% Generated draft: requires mathematical and source review.
\documentclass[11pt]{article}
\usepackage[margin=25mm]{geometry}
\usepackage{fontspec}
\setmainfont{Latin Modern Roman}
\usepackage{amsmath,amssymb,mathtools,unicode-math}
\setmathfont{Latin Modern Math}
\usepackage{enumitem,needspace,xurl,hyperref}
\setlength{\parindent}{0pt}
\setlength{\parskip}{5pt}
\setlength{\emergencystretch}{3em}
\newcommand{\N}{\mathbb{N}}
\newcommand{\Z}{\mathbb{Z}}
\newcommand{\Q}{\mathbb{Q}}
\newcommand{\R}{\mathbb{R}}
\newcommand{\C}{\mathbb{C}}
\newcommand{\F}{\mathbb{F}}
\newcommand{\A}{\mathbb{A}}
\newcommand{\PP}{\mathbb{P}}
\newcommand{\E}{\mathbb{E}}
\newcommand{\eps}{\varepsilon}
\newcommand{\dd}{\,\mathrm{d}}
\begin{document}
\section*{UnsolvedMath ID: 10400175}
\textbf{Source problem code:} catalogue entry
\subsection{Definitions and mathematical statement}
Let \(N\subset M\) be a subfactor for which the Turaev--Viro--Ocneanu invariant and the source's braiding are defined, and suppose the braiding is degenerate. In the usual finite semisimple ribbon-category interpretation, degeneracy means that the \(S\)-matrix is singular; equivalently, under the standard hypotheses, there are nontrivial transparent objects, whose double braiding with every object is the identity.

N. Sato asks for a description of the resulting Turaev--Viro--Ocneanu invariant as a Reshetikhin--Turaev invariant. More explicitly, the desired construction should assign suitable Reshetikhin--Turaev input data \(\mathcal D\) to the given subfactor and establish an identity of invariants
\[
 Z^{\mathrm{TVO}}_{N\subset M}(Y)=Z^{\mathrm{RT}}_{\mathcal D}(Y)
\]
for every \(3\)-manifold \(Y\) in their common intended domain, with the normalizations and any orientation, framing or other auxiliary data specified. The excerpt does not identify that manifold domain, the subfactor hypotheses or its braided category, or the required Reshetikhin--Turaev input. The displayed equality states the requested comparison, not a known identity or a proposed choice of \(\mathcal D\).
\subsection{Short English statement}
Express the Turaev--Viro--Ocneanu invariant of a subfactor with degenerate braiding as a Reshetikhin--Turaev invariant, including the required input category and normalization.
\subsection{Sources}
\begin{itemize}
\item N. Sato, attributed Problem 9.11 in T. Ohtsuki, \emph{Problems on invariants of knots and 3-manifolds} (2002), supplied PDF page 141. \url{https://msp.org/gtm/2002/04/gtm-2002-04-024s.pdf}. Recorded access 2026-07-29; extraction pdf-font-structure. No further citation is supplied for this item.
\item Local UnsolvedMath registry, \texttt{data/unsolvedmath/problems.json}, permanent ID \(10400175\), exact source code \texttt{catalogue entry}; research classification OPEN-TRIAGE.
\end{itemize}
\end{document}
